% Condensed Part 2 (Pretraining Language Models). Three reused result tables carry
% the SOTA claims (encoders, the CLM detour, OntoBook), plus the detour schema.
% All floats copied verbatim from sources/part_2/chapter{4,5,6}.

\noindent We then train compact biomedical encoders. Continuing a general French model's
pretraining on \emph{biomed-fr} improves all five biomedical NER benchmarks tested
(\Cref{tab:bioVSbase}) without training a new one from scratch, for a 2.54-point
average gain and at a fraction of the cost. CamemBERT-bio was adapted for an estimated
0.80~kg of CO\textsubscript{2}, against 26.11 for the from-scratch DrBERT.

\begin{table}[ht]
\centering
\small
\setlength{\tabcolsep}{5pt}
\sisetup{detect-weight=true}
\begin{threeparttable}[b]
\begin{tabular}{l c c c}
\toprule
& & \multicolumn{2}{c}{\textbf{CamemBERT-bio}} \\
\cmidrule(lr){3-4}
\textbf{Dataset} & \textbf{CamemBERT} & \textbf{biomed-fr-small} & \textbf{biomed-fr} \\
\midrule
\rowcolor{ThesisTableSep}
\multicolumn{4}{c}{\rule{0pt}{10pt}\textbf{Clinical}} \\[1pt]
CAS1 & $70.50 \pm 1.75$ & \underline{$72.94 \pm 1.12$} & $\mathbf{73.03 \pm 1.29}$ \\
CAS2 & $79.02 \pm 0.92$ & \underline{$80.00 \pm 0.32$} & $\mathbf{81.66 \pm 0.59}$ \\
E3C  & $67.63 \pm 1.45$ & \underline{$67.96 \pm 1.85$} & $\mathbf{69.85 \pm 1.58}$ \\
\midrule
\rowcolor{ThesisTableSep}
\multicolumn{4}{c}{\rule{0pt}{10pt}\textbf{Leaflets}} \\[1pt]
EMEA & $74.14 \pm 1.95$ & \underline{$75.93 \pm 2.42$} & $\mathbf{76.71 \pm 1.50}$ \\
\midrule
\rowcolor{ThesisTableSep}
\multicolumn{4}{c}{\rule{0pt}{10pt}\textbf{Scientific}} \\[1pt]
MEDLINE & \underline{$65.73 \pm 0.40$} & $65.48 \pm 0.31$ & $\mathbf{68.47 \pm 0.54}$ \\
\midrule
\textbf{Average} & 71.40 & \underline{72.46} & \bfseries 73.94 \\
\bottomrule
\end{tabular}
\begin{tablenotes}
  \small
  \item $F_1$ scores (micro-average, seqeval strict, IOB2). Mean $\pm$ standard deviation over 10 seeds. biomed-fr-s = biomed-fr-small (10\% subset).
\end{tablenotes}
\end{threeparttable}
\caption{\textbf{CamemBERT vs CamemBERT-bio} on five biomedical NER benchmarks. Continual pretraining on \textit{biomed-fr} yields +2.54 $F_1$ on average.}
\label{tab:bioVSbase}
\end{table}

We next introduce a causal detour: the model learns from next-token prediction,
then returns to standard bidirectional masked language modelling. As
\Cref{fig:clm-pipeline} shows, the detour keeps the architecture unchanged and only
swaps the attention mask and loss. During the causal phase the model cannot look
ahead, so it must compress information into each token representation, and the later
decay back to masking does not return it to an MLM-like state.

\begin{figure}[ht]
\centering
\begin{tikzpicture}[node distance=1.8cm, >=Stealth, scale=0.85, transform shape,
  box/.style={rectangle, rounded corners=3pt, draw=ThesisNeutral, minimum height=0.9cm, minimum width=2.4cm, align=center, font=\small\sffamily}
]
  \node[box, fill=ThesisPrimary!20] (base) {MLM encoder\\(ModernBERT)};
  \node[box, fill=ThesisTertiary!30, right=of base] (clm) {CLM phase\\(10--50B tokens)};
  \node[box, fill=ThesisSecondary!30, right=of clm] (decay) {MLM decay\\(10\% of Phase 1)};
  \node[box, fill=ThesisPrimary!30, right=of decay] (out) {ModernCamem-\\BERT-bio};

  \draw[->, thick, ThesisNeutral] (base) -- (clm) node[midway, above, font=\tiny\sffamily] {causal mask};
  \draw[->, thick, ThesisNeutral] (clm) -- (decay) node[midway, above, font=\tiny\sffamily] {bidir.\ mask};
  \draw[->, thick, ThesisNeutral] (decay) -- (out);
\end{tikzpicture}
\caption{\textbf{The CLM detour pipeline.} Starting from a pretrained MLM encoder, we continue pretraining with a causal objective (Phase~1), then decay back to MLM (Phase~2). The model does not return to an MLM-like state.}
\label{fig:clm-pipeline}
\end{figure}

On the eight French biomedical tasks the detour beats the matched MLM baseline, as
\Cref{tab:french-results} reports.

\begin{table}[ht]
\centering
\small
\caption{\textbf{French biomedical results} (macro $F_1$, 9 seeds). Bold: best per column; underline: second best.}
\label{tab:french-results}
\resizebox{\textwidth}{!}{
\begin{tabular}{lcccccccccc}
\toprule
& & \multicolumn{5}{c}{\textbf{Multilabel Classification}} & {\textbf{CLS}} & \multicolumn{2}{c}{\textbf{NER}} & \\
\cmidrule(lr){3-7} \cmidrule(lr){8-8} \cmidrule(lr){9-10}
& \textbf{Ctx} & \textbf{FRACCO-30} & \textbf{FRACCO-100} & \textbf{CANTEMIST} & \textbf{DISTEMIST} & \textbf{MedDialog} & \textbf{DiaMED} & \textbf{EMEA} & \textbf{Medline} & \textbf{Avg} \\
\midrule
\rowcolor{ThesisTableSep} \multicolumn{11}{c}{\rule{0pt}{10pt}\textbf{Baselines}} \\[1pt]
ModernCamemBERT & 8192 & 70.1 & 55.3 & 63.3 & 20.2 & 60.6 & 56.4 & 63.4 & 55.3 & 55.6 \\
DrBERT & 512 & 53.0 & 35.6 & 37.9 & 21.4 & \underline{63.6} & 57.0 & \textbf{68.0} & \textbf{62.3} & 49.9 \\
CamemBERT-bio & 512 & 41.9 & 20.1 & 12.8 & 9.6 & 38.6 & 47.7 & 61.6 & 56.6 & 36.1 \\
CamemBERT & 512 & 40.8 & 19.4 & 11.9 & 9.5 & 37.4 & 40.6 & 61.5 & 54.7 & 34.5 \\
\midrule
\rowcolor{ThesisTableSep} \multicolumn{11}{c}{\rule{0pt}{10pt}\textbf{Our Models (Base, 150M)}} \\[1pt]
MLM baseline & 8192 & 69.9 & 56.8 & 64.9 & 23.5 & 62.5 & 63.4 & 65.4 & 56.8 & 57.9 \\
CLM detour & 8192 & 74.8 & 60.1 & 71.0 & 25.5 & \underline{63.6} & \textbf{67.4} & 65.9 & 58.2 & 60.8 \\
\midrule
\rowcolor{ThesisTableSep} \multicolumn{11}{c}{\rule{0pt}{10pt}\textbf{Our Models (Large, 350M)}} \\[1pt]
MLM baseline & 8192 & \underline{79.4} & \underline{63.3} & \underline{72.6} & \underline{29.1} & \textbf{64.5} & 61.2 & 66.1 & 58.0 & \underline{61.8} \\
CLM detour & 8192 & \textbf{80.7} & \textbf{65.4} & \textbf{74.4} & \textbf{30.4} & \textbf{64.5} & \underline{64.8} & \underline{67.0} & \underline{59.5} & \textbf{63.3} \\
\bottomrule
\end{tabular}
}
\end{table}

It reaches 60.8\% average $F_1$ against 57.9\% at Base, and 63.3\% against 61.8\% at
Large. Since the two trajectories differ only in the Phase~1 objective, the gap is
the effect of that objective alone, at no extra cost, and it transfers to English.
We release these models as ModernCamemBERT-bio.

We also introduce \emph{OntoBook}, which complements text pretraining by converting
medical ontologies into synthetic textbooks: an ontology subgraph is turned into a
random walk and reformulated into fluent prose, and the encoder trains on it with
two aligned objectives, masked language modelling and relation prediction between
code pairs. It outperforms all baselines on the three French coding benchmarks, with
the largest gain on the fine-grained DISTEMIST coding, where it also cuts the
variance from $\pm$5.8 to $\pm$1.1.

We release the resulting models and data.
